% BHU PYQ -- Professional Exam Template
\documentclass[11pt,a4paper]{article}

\usepackage[a4paper,top=2.5cm,bottom=2.5cm,left=2.8cm,right=2.8cm,headheight=14pt]{geometry}
\usepackage{amsmath,amssymb}
\usepackage[T1]{fontenc}
\usepackage[utf8]{inputenc}
\usepackage{lmodern}
\usepackage{microtype}
\usepackage{fancyhdr}
\usepackage{enumitem}
\usepackage{hyperref}
\usepackage{lastpage}
\usepackage{parskip}

\hypersetup{colorlinks=true,urlcolor=black,linkcolor=black,
  pdftitle={B.Sc. (Hons.) Botany Sem-VI BOB-603 -- BHU PYQ}}

\pagestyle{fancy}
\fancyhf{}
\renewcommand{\headrulewidth}{0.4pt}
\renewcommand{\footrulewidth}{0.4pt}
\fancyhead[L]{\small   \textit{Banaras Hindu University}}
\fancyhead[R]{\small   \textit{Botany Paper: BOB-603}}
\fancyfoot[C]{\small Page  \thepage\ of \pageref{LastPage}}


\newlist{parts}{enumerate}{2}
\setlist[parts,1]{label=(\alph*),leftmargin=*,itemsep=4pt,topsep=3pt}
\setlist[parts,2]{label=(\roman*),leftmargin=*,itemsep=2pt,topsep=2pt}

\newcommand{\pts}[1]{\hfill   \textit{[#1]}}


\begin{document}


\begin{center}
  {\large\bfseries BANARAS HINDU UNIVERSITY}\\[2pt]
  {\normalsize B.Sc. (Hons.) Semester VI Examination, 2022-23}\\[6pt]
  \rule{0.8\linewidth}{0.5pt}\\[6pt]
  {\large\bfseries Botany}\\[4pt]
  {\large\bfseries Paper: BOB-603 --- Cytogenetics and Evolutionary Processes}\\[4pt]
  {\normalsize Time: Three hours \hfill Maximum Marks: 70}
\end{center}

\rule{\linewidth}{0.4pt}

\smallskip



\noindent   \textit{Instructions: Answer five questions in all, selecting two each from Sections A and B and Section-C which is compulsory. The figures in the right-hand margin indicate marks.}

\bigskip


\begin{center}
   \textbf{SECTION A}
\end{center}
\medskip

\medskip


\noindent   \textbf{Question 1.} Describe the various types of chromosomal aberrations with suitable examples.\pts{15}

\medskip


\noindent   \textbf{Question 2.} Describe any two of the following:\pts{$2  \times 7.5 = 15$}

\begin{parts}
  \item Lampbrush chromosome
  \item Chemical composition of chromosome
  \item Difference between inversion and translocation with examples
\end{parts}

\medskip


\noindent   \textbf{Question 3.} What is karyotype? How can it be used to describe evolution? Illustrate with suitable examples.\pts{15}

\medskip


\noindent   \textbf{Question 4.} Write critical notes on any three of the following:\pts{$3  \times 5 = 15$}

\begin{parts}
  \item Linkage map
  \item Chromatin and its types
  \item Heteroploidy
  \item Different types of trisomic
\end{parts}

\bigskip

\begin{center}
   \textbf{SECTION B}
\end{center}
\medskip

\medskip


\noindent   \textbf{Question 5.} Explain how recombination frequency can be used to determine gene positions on a chromosome. Illustrate with examples.\pts{15}

\medskip


\noindent   \textbf{Question 6.} Explain any two of the following:\pts{$2  \times 7.5 = 15$}

\begin{parts}
  \item Deletion mapping
  \item Somatic cell fusion and selection of hybrid cells
  \item Difference between alkylating and intercalating agents
\end{parts}

\medskip


\noindent   \textbf{Question 7.} What is cytoplasmic inheritance? Describe cytoplasmic inheritance with suitable examples.\pts{15}

\medskip


\noindent   \textbf{Question 8.} Discuss any two of the following:\pts{$2  \times 7.5 = 15$}

\begin{parts}
  \item Difference between qualitative and quantitative trait inheritance
  \item Multiple allelism
  \item Mitochondrial and chloroplast DNA
\end{parts}

\bigskip

\begin{center}
   \textbf{SECTION C}
\end{center}
\medskip

\medskip


\noindent   \textbf{Question 9.} Write short notes on the following in not more than 50 words each:\pts{$10  \times 1 = 10$}

\begin{parts}
  \item Base analogs
  \item 6,4-PP
  \item Radiation hybrids
  \item GISH
  \item Reciprocal cross
  \item   \textit{Brassicoraphanus}
  \item Centromere
  \item Eukaryotic organelles involved in cytoplasmic inheritance
  \item Biochemical nature of Rh factor
  \item Self-incompatibility.
\end{parts}

\vfill
\rule{\linewidth}{0.4pt}

\begin{center}\small
\href{https://bhu-exam.vercel.app/}{Click Here}\\[4pt]\href{https://me-aryan.vercel.app/}{Click Here}\\[4pt]\href{https://quashan-me.vercel.app/}{Click Here}
\end{center}

\end{document}
